\section{Fonctionnalités avancées et industrialisation}
\label{sec:avance}

Ce chapitre présente les évolutions qui portent la plateforme au-delà du socle
initial : documentation d'API, tests de charge, observabilité complète (les trois
piliers métriques / journaux / traces), alerting orienté SLO, cycle de vie et
qualité des données, chaîne CI/CD sécurisée et déploiement Kubernetes.

\subsection{Documentation d'API (OpenAPI / Swagger)}
L'API REST est désormais auto-documentée à l'aide de \texttt{drf-spectacular}, qui
génère un schéma OpenAPI~3 à partir des vues et des sérialiseurs
existants~\cite{drf-spectacular}. Trois points d'entrée sont exposés :
\texttt{/api/schema/} (schéma brut), \texttt{/api/docs/} (interface Swagger UI) et
\texttt{/api/redoc/} (ReDoc). Chaque endpoint applicatif est annoté afin d'obtenir
un schéma sans erreur.

\begin{figure}[H]
    \centering
    \includegraphics[width=0.9\textwidth]{swagger_ui.png}
    \caption{Documentation interactive de l'API (Swagger UI)}
    \label{fig:swagger-ui}
\end{figure}

\subsection{Tests de charge (k6)}
Une suite \texttt{k6} (\path{loadtest/ingest_and_read.js}) génère une montée en
charge progressive contre les endpoints d'ingestion et de lecture, en injectant
une proportion configurable d'événements en erreur (5xx) afin de faire réagir la
chaîne d'observabilité~\cite{k6-docs}. Des seuils (\emph{thresholds}) transforment
le test en garde-fou de performance.

\begin{verbatim}
make demo        # démarrer la stack
make loadtest    # BASE_URL / BATCH / ERROR_RATIO paramétrables
\end{verbatim}

\begin{figure}[H]
    \centering
    \includegraphics[width=0.9\textwidth]{k6_summary.png}
    \caption{Résultats du test de charge k6 (latence p95, taux d'erreur)}
    \label{fig:k6-summary}
\end{figure}

\subsection{Les trois piliers de l'observabilité}
La plateforme couvre désormais les trois piliers de l'observabilité, tous
consultables dans Grafana.

\subsubsection{Métriques}
Prometheus collecte les métriques applicatives (\texttt{django-prometheus} et des
métriques métier telles que \texttt{apm\_ingested\_requests\_total} et
\texttt{apm\_ingest\_latency\_seconds}), ainsi que \texttt{node-exporter} et
\texttt{postgres-exporter}.

\subsubsection{Traces distribuées (OpenTelemetry + Tempo)}
L'application Django est instrumentée avec OpenTelemetry ; les spans sont exportés
en OTLP vers un OpenTelemetry Collector qui les transmet à Grafana
Tempo~\cite{opentelemetry-docs,grafana-tempo}. Le traçage est activable via la
variable \texttt{OTEL\_ENABLED} et reste désactivé pendant les tests et
l'intégration continue.

\begin{figure}[H]
    \centering
    \includegraphics[width=0.9\textwidth]{tempo_trace.png}
    \caption{Trace distribuée d'une requête (Grafana Tempo)}
    \label{fig:tempo-trace}
\end{figure}

\subsubsection{Journaux structurés (Loki)}
L'application émet des journaux au format JSON incluant \texttt{trace\_id} et
\texttt{span\_id}. Promtail collecte les journaux des conteneurs et les expédie
vers Loki~\cite{grafana-loki}. Grafana relie les traces aux journaux (Tempo
$\rightarrow$ Loki) et inversement (Loki $\rightarrow$ Tempo).

\begin{figure}[H]
    \centering
    \includegraphics[width=0.9\textwidth]{loki_logs.png}
    \caption{Journaux structurés corrélés aux traces (Grafana Loki)}
    \label{fig:loki-logs}
\end{figure}

\subsection{Alerting orienté SLO}
Au-delà des alertes de seuil (taux d'erreur, latence, CPU), la plateforme définit
un objectif de niveau de service (SLO) de disponibilité à 99,5\,\% et des alertes
de consommation du budget d'erreur (\emph{burn-rate}) multi-fenêtres :
combustion rapide (14,4$\times$ sur 1\,h et 5\,min $\rightarrow$ alerte critique) et
combustion lente (6$\times$ sur 6\,h et 30\,min $\rightarrow$ avertissement), selon la
méthodologie du \emph{Google SRE Workbook}~\cite{google-sre-slo}.

\begin{figure}[H]
    \centering
    \includegraphics[width=0.9\textwidth]{prometheus_alerts.png}
    \caption{Alertes de budget d'erreur (Prometheus / Alertmanager)}
    \label{fig:prometheus-alerts}
\end{figure}

\subsection{Cycle de vie et qualité des données}
La migration \texttt{0009\_retention\_compression} active la compression
colonnaire native de TimescaleDB (par défaut au-delà de
\texttt{APM\_COMPRESS\_AFTER\_DAYS} = 7 jours) ainsi qu'une politique de rétention
(\texttt{APM\_RETENTION\_DAYS} = 90 jours) sur l'hypertable brute, les vues
continues conservant l'historique agrégé de long terme~\cite{timescaledb-compression}.

Une commande \texttt{check\_data\_quality} vérifie l'absence de valeurs manquantes,
la validité des plages (\texttt{status\_code}, \texttt{latency\_ms}), le taux de
\texttt{trace\_id} dupliqués et la fraîcheur des données. Elle renvoie un code de
sortie non nul en cas de violation, ce qui permet de l'utiliser comme garde-fou en
intégration continue.

\begin{verbatim}
python manage.py check_data_quality --max-age-minutes 60 --fail-on-empty
\end{verbatim}

\subsection{Chaîne CI/CD sécurisée}
Le pipeline d'intégration continue (lint, formatage, tests SQLite et TimescaleDB,
smoke Docker Compose, audit des dépendances) est complété par un garde-fou de
qualité des données et par des contrôles de sécurité : analyse statique CodeQL et
scan de vulnérabilités de l'image conteneur avec Trivy~\cite{trivy-docs}. La
livraison continue construit et pousse l'image vers GHCR avec génération d'un SBOM,
attestation de provenance et signature Cosign.

\subsection{Déploiement Kubernetes (Helm + GitOps)}
En complément de Docker Compose et d'Ansible, un chart Helm
(\path{deploy/helm/apm-observability}) décrit le déploiement de l'application sur
Kubernetes (Deployment, Service, Ingress, ConfigMap/Secret, HPA, et une
TimescaleDB en cluster optionnelle)~\cite{helm-docs}. Une \emph{Application} Argo CD
(\path{deploy/argocd/application.yaml}) assure une livraison déclarative en GitOps :
le cluster est maintenu en synchronisation avec le chart versionné dans Git, avec
auto-réparation et élagage des ressources~\cite{argocd-docs}.

\begin{verbatim}
helm lint deploy/helm/apm-observability
helm install apm deploy/helm/apm-observability --namespace apm --create-namespace
kubectl apply -f deploy/argocd/application.yaml   # variante GitOps
\end{verbatim}

\begin{figure}[H]
    \centering
    \includegraphics[width=0.9\textwidth]{argocd_app.png}
    \caption{Livraison GitOps de la plateforme (Argo CD)}
    \label{fig:argocd-app}
\end{figure}
